\section{Width bounds and the cost of the transfer}
\label{sec:width}

\begin{lemma}[LOW width]
\label{lem:lowwidth}
\lean{Meanders.FirstCrossing.low_geometry_port_bound, Meanders.FirstCrossing.high_geometry_port_bound}
\leanok
\uses{def:sectors, def:key}
In \LOW\ the number of retained ports after any point is
$h_P+h_Q=2k\le2(K-1)$.
\end{lemma}

\begin{proof}
\leanok
With the chronological schedule there is no right side and no emission; the
ports are the unmatched openings, $h_P+h_Q=2k$ of them, and $k\le K-1$ is
enforced after every point.
\end{proof}

\begin{lemma}[HIGH width]
\label{lem:highwidth}
\lean{Meanders.FirstCrossing.high_geometry_port_bound}
\leanok
\uses{lem:partition, lem:forced, def:key}
In a \HIGH\ sector with $d=n-K$, the number of retained ports satisfies
\begin{center}
\begin{tabular}{ll}
\toprule
stage of the schedule & retained ports \\
\midrule
after $t\le\Delta$ initial long-side points & $\le2t\le2d$ \\
after a completed long/short pair & $\le2d$ \\
between the two points of a pair & $\le2d+2$ \\
with $t\le\Delta$ final long-side points unread & $\le2t\le2d$ \\
\bottomrule
\end{tabular}
\end{center}
A single transition may temporarily hold $2d+4$ ports before its forced joins.
\end{lemma}

\begin{proof}[Proof sketch]
\leanok
Initial skirt: each processed point contributes at most two ports and nothing
has been joined, so $p\le2t$, and $t\le\Delta\le d$ by \cref{lem:partition}.
Completed pair: for one owner with cut height $w$, the identity
$p=A_L+A_R-2c_{\min}-2\max(0,c_{\max}-I)$ of \cref{app:ports}, where
$I=i_L-A_L=i_R-A_R$, $c_{\max}=\max(c_L,c_R)$ and $c_{\min}=\min(c_L,c_R)$,
gives $p\le A_L+A_R=n-w$; summing over the two owners gives $2n-u-v=2d$. Final skirt: with $t$ long points unread and
$r$ of the long budget remaining, $p=t-2\max(0,t-w-r)\le t$ per owner
(\cref{app:ports}). The middle row adds the two ports of one point to the
completed-pair bound.
\end{proof}

\begin{proposition}[Carrier size]
\label{prop:carrier}
\lean{Meanders.FirstCrossing.carrierCard_le, Meanders.FirstCrossing.low_carrierCard_le, Meanders.FirstCrossing.high_carrierCard_le}
\leanok
\uses{lem:lowwidth, lem:highwidth, lem:seams}
The number of distinct keys in any layer is at most
\[
 \sum_{k=0}^{K-1}(2k+1)\Cat_k=O\bigl(n\,4^{K-1}\bigr)\ \text{in \LOW},\qquad
 (2n+1)^{4}\sum_{j=0}^{d+1}\Cat_j=O\bigl(n^{4}4^{\,d}\bigr)\ \text{in a \HIGH\ sector}.
\]
\end{proposition}

\begin{proof}
\leanok
In \LOW\ the key is determined by the split of $2k$ ports between $P$ and $Q$
($2k+1$ possibilities) and a noncrossing pairing of $2k$ positions. In \HIGH\
the four counters take at most $(2n+1)^{4}$ values, and the pairing is a
noncrossing pairing of at most $2d+2$ positions.
\end{proof}

The Lean statements corresponding to this proposition are
$(n+1)^{4}\,K\,4^{K-1}$ for \LOW\ and $(n+1)^{4}\sum_{j=0}^{d+1}\Cat_j$ for a
\HIGH\ sector, both for the set of keys that pass the validity test of a
stage; they carry the same exponential factors with looser polynomial
prefactors than the display above, and the paper's sharper \LOW\ form is not
formalized (\cref{sec:evidence}).

\begin{theorem}[Cost for a fixed threshold]
\label{thm:cost}
\uses{lem:partition, lem:future, lem:cycle, prop:carrier, lem:lowwidth, lem:highwidth}
For $1\le K\le n$ the transfer computes $\Closed(n)$ in
$\operatorname{poly}(n)\bigl(4^{K-1}+4^{n-K}\bigr)$ bit time and bit space.
\end{theorem}

\begin{proof}
There are $O(n^{2})$ sectors, each with $O(n)$ stages, processed sequentially
with two logical live layers plus the working storage of one sort. By \cref{prop:carrier} a layer has $O(n^{4}4^{n-K})$
keys in \HIGH\ and $O(n4^{K-1})$ in \LOW, each of $O(n\log n)$ bits. Each
key emits at most four successors; one trial performs a constant number of
adjacent surgeries, at most two forced joins, and a canonical renaming, all in
polynomial time on arrays of length $O(n)$. If a layer emits $M$ records,
deterministic comparison sorting of the records by key costs $O(M\log M)$
comparisons of $O(n\log n)$-bit keys, and $\log M=O(n)$ because $M$ is at
most four times the key bound; merging equal keys is then one linear pass.
The weight $C$ of a key after $t$ points counts surviving labelled
histories of length $t$, so $C\le4^{t}\le4^{2n}$ has $O(n)$ bits; the
observer weight of \cref{sec:open} satisfies $0\le E\le nC$, since a history
has at most $n$ lower $D$ letters, so it also has $O(n)$ bits, and merging
adds such integers. Sorting the $M$ emitted records needs $O(M)$ records of
workspace, within the same bound. Summing over stages and sectors gives the
bound, with the two exponentials arising from the two branches.
\end{proof}

\begin{corollary}[Balance]
\label{cor:balance}
\uses{thm:cost}
With $K=\lfloor n/2\rfloor+1$ the transfer computes $\Closed(n)$ in
$\Ostar(2^{n})$ bit time and space.
\end{corollary}

\begin{proof}
For $n=2m$ we have $K=m+1$, $d=m-1$, so $4^{K-1}=2^{n}$ and $4^{d}=2^{n-2}$.
For $n=2m+1$ we have $K=m+1$, $d=m$, so $4^{K-1}=4^{d}=2^{n-1}$.
\end{proof}

\begin{remark}
The count of keys per layer is an upper bound on what is stored; the
transfer generates only keys reachable from the empty seed, so its actual
memory is often much smaller. No enumeration of complete half words, no
precomputed table and no completion oracle is used: the Catalan bounds are
never traversed as lists.
\end{remark}
